\textbf{(a) A larger smallest clean eigenvalue improves recovery; (b) more noise increases instability; (c) more pairs reduce subspace error near $k^{-1/2}$.} GRIT ($r=2$) on the Gaussian synthetic task. Circles show subspace error $\|(I-P_{k,2})P_{\mathcal U}\|_{\rm op}$; squares show normalized instability $\|P_{\mathcal U}\boldsymbol w\|_2/\|\boldsymbol w\|_2$, bounded above by that error for each probe. For $k\ge512$, mean log--log slopes are -0.497 and -0.864, respectively; the faster predictor decay is empirical, not implied by the bound. All panels fix 96 features, 3,000 labels, and ridge $0.1$. Labeled data and target distributions are shared across conditions. Pair differences receive independent isotropic Gaussian noise; $\sigma_\xi$ is its coordinate standard deviation. Panel (a) fixes $k=256$, $\sigma_\xi=0.1$, and $\lambda_{\max}=1$, so clean trace varies. Panels (b,c) fix clean spectrum $(1,1)$; (b) fixes $k=256$ and (c) fixes $\sigma_\xi=0.1$. Lines show means and bands show sample SD over 20 replicates. Vertical axes are linear in (a,b), logarithmic in (c). The dotted shape guide is anchored to mean subspace error at prespecified $k=256$, not calibrated as a bound. Zero-noise and nonspanning controls are reported separately. 
